\section{Supplementary verification notes}\label{app:reproducibility-note}

The proof-bearing computations are the algebraic identities displayed in the article: the Sylvester determinant in Appendix~\ref{app:disc}, the cubic branch-resultant determinant in Lemma~\ref{lem:cubic-branch-elimination}, the positive-axis endpoint and opposite-modulus resultants in Proposition~\ref{prop:finite-interval-certificates}, the numerator--denominator resultant and Sturm pseudo-remainders in Appendix~\ref{app:degeneracies}, and the endpoint expansion package in Appendix~\ref{app:endpoint-branching}.
Appendix~\ref{app:projective-infinity} has a different status: it supplies a supplementary large-\(y\) consistency check and is not used in the endpoint proof.
The following numerical comparisons are independent checks of the same algebraic conclusions, not hypotheses in any proof.
\begin{enumerate}
\item The real-axis ordering of the moduli of the four roots of \(\Pi(\lambda,y)=0\) agrees with Theorem~\ref{thm:ep-classification} and Proposition~\ref{prop:global-dominance}: the branch value \(y=0\) is the only maximal-modulus collision on the positive side, while \(y=1\) and \(y=\ySub\) are non-dominant collisions.
\item The first negative zeros obtained from the recurrence for \(m=80,120,200\) have scaled values \(-m^2y_{m,k}\), for fixed \(k\), matching the leading constants \(\pi^2(2k-1)^2/2\) and the two-term correction in Theorem~\ref{thm:endpoint-zero-quantization}.
\item The factor \((y+1)^2\) from Proposition~\ref{prop:double-zero} appears in the tested recurrence polynomials for \(m\ge2\), while no corresponding forced zero appears at the non-dominant branch value \(y=1\).
\end{enumerate}

The script \path{scripts/generate_validation_figure.py} reconstructs \(Z_m(y)\) from the recurrence \eqref{eq:recurrence}, evaluates the four roots of \eqref{eq:spectral-poly} near the real branch values, removes the algebraically forced factor \((y+1)^2\) from Proposition~\ref{prop:double-zero}, and compares the first negative zeros with Theorem~\ref{thm:endpoint-zero-quantization}.
All logical dependencies needed for the theorems are contained in the labelled results and displayed identities of the article itself.
